
\blank
\newpage

\section{Introduction}

    \subsection{Cos'è il Natural Language Processing}
    Il \textbf{Natural Language Processing} (NLP), o Elaborazione del Linguaggio Naturale, è l'area
    dell'Intelligenza Artificiale che si occupa dello studio e dello sviluppo di metodi e sistemi
    in grado di analizzare, comprendere e generare linguaggio umano.
    A differenza dei linguaggi formali o di programmazione, il linguaggio naturale è intrinsecamente ambiguo,
    ricco di sfumature semantiche e fortemente legato al contesto.  
    L'NLP rappresenta dunque una disciplina di frontiera, che integra competenze linguistiche, informatiche e statistiche
    per trasformare testi e conversazioni in conoscenza computabile.
    
    \subsection{Obiettivi del corso}
    L'obiettivo del corso è fornire una panoramica aggiornata delle più recenti tecnologie di
    Intelligenza Artificiale applicate al linguaggio. In particolare, si intende mostrare come un robot
    o un agente virtuale possano acquisire la capacità di comprendere ed elaborare il linguaggio umano
    per diversi scopi applicativi.
    
    \subsection{Competenze acquisite}
    Alla fine del corso, lo studente sarà in grado di comprendere e utilizzare strumenti che permettono di:
    
    \begin{itemize}
      \item \textbf{Rispondere a domande} tramite modelli di \emph{question answering}, capaci di cercare e sintetizzare informazioni in modo coerente.
      \item \textbf{Riconoscere la struttura di un testo}, ad esempio individuando frasi, dipendenze sintattiche e categorie grammaticali.
      \item \textbf{Individuare tratti emotivi e di genere} all'interno di un documento, sfruttando tecniche di \emph{sentiment analysis} e classificazione.
      \item \textbf{Tradurre tra lingue diverse}, mediante modelli neurali di traduzione automatica basati su architetture Transformer.
      \item \textbf{Interagire in modo fluente con gli utenti}, integrando modelli di dialogo e assistenti conversazionali.
    \end{itemize}
    
    \subsection{Prospettive applicative}
    Queste competenze trovano applicazione in numerosi contesti: dai sistemi di assistenza virtuale
    all'analisi automatica di grandi collezioni di testi, fino alla creazione di strumenti multilingue
    per la comunicazione e la ricerca.  
    Il corso costituisce dunque una base solida per chi desidera approfondire il ruolo delle tecnologie linguistiche nella società digitale contemporanea.



    